% s-110-p110-m003; original PDF 110, 111
\Aphorism{s-110-p110-m003}{119}
Wo die Fehltritte der Weiber von Recht und Sitte hart beurtheilt werden, wie im alten Rom, da ist wahrscheinlich Sitte und Recht in den Händen der Besitzenden, der Ehemänner und Familienväter, oder der Besitzsicheren, die ein Weib zu erwerben und festzuhalten wissen; Untreue gilt hier als Vertragsbruch, als Empörung gegen die herrschende Kaste. Wo umgekehrt dem ausschweifenden Weibe milde begegnet wird, da führt der Instinct der „unglücklichen Liebhaber“, der erotisch Declassirten das Wort, mit dem Hintergedanken: Wenn es nicht solche Weiber gäbe, wie kämen \Sperr{wir} je zu einem Schäferstündchen? Es ist das Schwächegefühl gegenüber stolzen und ehrbaren Frauen, das keine Strenge gegen die verdorbenen aufkommen lässt.
\par

